% !TEX program = pdflatex
% CV français ciblé pour le dossier RH Inria Startup Studio / Percolia.
% Compilation : pdflatex -interaction=nonstopmode -halt-on-error CV_Louis_Hauseux_Inria_Startup_Studio_2026.tex
\documentclass[10pt,a4paper]{article}

\usepackage[T1]{fontenc}
\usepackage[utf8]{inputenc}
\usepackage{helvet}
\renewcommand{\familydefault}{\sfdefault}
\usepackage[left=1.35cm,right=1.35cm,top=0.95cm,bottom=0.95cm]{geometry}
\usepackage{microtype}
\usepackage{xcolor}
\usepackage{titlesec}
\usepackage{enumitem}
\usepackage{tabularx}
\usepackage{ragged2e}
\usepackage{hyperref}

\definecolor{darkblue}{RGB}{20,52,80}
\definecolor{midblue}{RGB}{53,94,128}
\definecolor{softgray}{RGB}{78,78,78}

\hypersetup{
  unicode=true,
  colorlinks=true,
  urlcolor=darkblue,
  linkcolor=darkblue,
  pdfauthor={Louis Hauseux},
  pdftitle={Curriculum vitae -- Louis Hauseux -- Inria Startup Studio},
  pdfsubject={Dossier RH Inria Startup Studio -- projet Percolia},
  pdfkeywords={Inria Startup Studio, Percolia, HGP-Clusterer, clustering hiérarchique, LiDAR 3D, géométrie algorithmique}
}
\urlstyle{same}

\titleformat{\section}
  {\normalsize\bfseries\color{darkblue}}
  {}{0pt}{\MakeUppercase}[\vspace{-0.30em}\color{midblue}\rule{\textwidth}{0.45pt}]
\titlespacing*{\section}{0pt}{0.68em}{0.34em}

\setlength{\parindent}{0pt}
\setlength{\tabcolsep}{0pt}
\setlength{\emergencystretch}{2em}
\newcolumntype{L}[1]{>{\RaggedRight\arraybackslash\hyphenpenalty=10000\exhyphenpenalty=10000}p{#1}}
\newcolumntype{Y}{>{\RaggedRight\arraybackslash\hyphenpenalty=10000\exhyphenpenalty=10000}X}
\setlist[itemize]{
  leftmargin=3.60cm,
  label=\textcolor{midblue}{\raisebox{0.12ex}{\scriptsize$\bullet$}},
  labelsep=0.42em,
  itemsep=0.04em,
  topsep=0.05em,
  parsep=0pt,
  partopsep=0pt
}

% #1 période ; #2 durée/statut ; #3 employeur/projet ; #4 lieu ; #5 fonction.
\newcommand{\cventry}[5]{%
  \begin{tabularx}{\textwidth}{@{}L{3.23cm}@{\hspace{0.25cm}}Y@{}}
    \textcolor{softgray}{#1} & \textbf{#3}\hfill\textcolor{softgray}{#4}\\[-0.10ex]
    \small\textcolor{darkblue}{\textbf{#2}} & \textbf{#5}
  \end{tabularx}\vspace{-0.20em}
}

\newcommand{\compactentry}[2]{%
  \begin{tabularx}{\textwidth}{@{}L{3.23cm}@{\hspace{0.25cm}}Y@{}}
    \textcolor{softgray}{#1} & #2
  \end{tabularx}\par
}

\begin{document}
\pagestyle{empty}
\frenchspacing
\RaggedRight
\hyphenpenalty=10000
\exhyphenpenalty=10000

% ==================== EN-TÊTE ====================
\begin{center}
  {\LARGE\bfseries Louis HAUSEUX}\\[0.20em]
  {\large\bfseries\color{darkblue}Porteur scientifique et technologique de Percolia}\\[0.10em]
  {\normalsize Normalien (ENS Paris-Saclay) \textbullet{} Doctorant en informatique -- soutenance le 8 septembre 2026}\\[0.28em]
  Sophia Antipolis
  \enspace\textcolor{midblue}{\textbullet}\enspace
  \href{mailto:louis.hauseux@inria.fr}{louis.hauseux@inria.fr}
  \enspace\textcolor{midblue}{\textbullet}\enspace
  \href{https://github.com/Ludwig-H}{GitHub}
  \enspace\textcolor{midblue}{\textbullet}\enspace
  \href{https://inria.hal.science/search/index/?q=authIdHal_s:louis-hauseux}{HAL}
\end{center}
\vspace{-0.48em}

% ==================== PROFIL ====================
\section*{Profil}
Normalien, mathématicien et informaticien, spécialisé en structuration automatique de
données, clustering hiérarchique, géométrie algorithmique et traitement de nuages de points
3D. Au cours de mon doctorat, j'ai conçu \textbf{HGP-Clusterer}, noyau scientifique du projet
\textbf{Percolia}. J'en porte la R\&D, la robustesse, le passage à l'échelle et la transformation
en produit logiciel. J'interviens en complémentarité avec Alban Hauseux, porteur business,
chargé de la découverte client, des audits comparatifs et de l'intégration.

% ==================== PERCOLIA ====================
\section*{Projet entrepreneurial -- Percolia}
\cventry{2025--2026}{Entrée ISS : 01/10/2026}{Projet Percolia -- Inria Startup Studio}{Sophia Antipolis}{Porteur technologique -- R\&D, passage à l'échelle et mise en produit}
\begin{itemize}
  \item Concevoir et prototyper HGP-Clusterer, méthode de clustering fondée sur les
  interactions d'ordre supérieur et la percolation pour limiter les connexions parasites dues
  au bruit.
  \item Structurer des nuages LiDAR 3D sans apprentissage ni données annotées, avec injection
  possible d'a priori géométriques ; prototype éprouvé sur un benchmark de 2 millions de
  points et sur des données réelles de l'ordre de 10 millions de points.
  \item Pendant ISS, industrialiser le noyau, établir les benchmarks clients et préparer son
  déploiement sous forme d'API ou de licence locale, ainsi que la stratégie de propriété
  intellectuelle.
\end{itemize}

\smallskip
{\normalsize\bfseries\color{darkblue}Expertise mise au service de Percolia}\par
\smallskip
\compactentry{R\&D algorithmique}{Hiérarchies de densité, graphes et hypergraphes,
percolation et géométrie algorithmique ; fondements publiés dans \textit{Applied Network
Science} en 2026.}
\compactentry{Logiciel / données}{Python, PyTorch/GPU, R et Git ; prototypage, optimisation
d'algorithmes et traitement de données massives.}
\compactentry{LiDAR / validation}{Détection d'anomalies 3D guidée par un modèle théorique ;
segmentation et suivi d'instances 4D sur SemanticKITTI, sans apprentissage.}
\compactentry{Transfert / partenaires}{Contrats Inria avec Airbus Defence and Space et Nokia ;
collaboration Cerema ; formation Starthèse et actions Inria Startup Studio.}

% ==================== EXPÉRIENCE PAGE 1 ====================
\section*{Expériences professionnelles}
\cventry{Oct. 2023--sept. 2026}{36 mois}{Université Côte d'Azur}{Sophia Antipolis}{Doctorant contractuel en informatique -- accueilli au Centre Inria, équipes NEO et AYANA}
\begin{itemize}
  \item Définir, analyser et implémenter HGP-Clusterer : correspondance exacte avec la
  hiérarchie des amas de forte densité $K$-NN et comparaison à Robust Single-Linkage,
  DBSCAN et HDBSCAN par la percolation.
  \item Décliner ces méthodes en LiDAR 3D/4D, traitement d'images multimodales et graphes
  pondérés signés ; produire des démonstrateurs, publications et collaborations appliquées.
  \item Assurer, en mission concomitante, \textbf{154,125 hETD} en M1/M2 Data Science and
  Artificial Intelligence : analyse de graphes, inférence statistique et apprentissage
  statistique.
\end{itemize}

\cventry{Fév.--sept. 2023}{8 mois}{Inria}{Sophia Antipolis}{Ingénieur de recherche contractuel}
\begin{itemize}
  \item Équipe AYANA, 3 mois, contrat Airbus Defence and Space : télédétection, détection et
  suivi d'objets sur images satellitaires, synthèse et valorisation scientifiques.
  \item Équipe NEO, 5 mois, contrat Nokia : estimation de densité, graphes géométriques et
  clustering ; développement d'un prototype et contributions scientifiques.
\end{itemize}

\cventry{2021--2022}{Année scolaire}{Lycée EMC2--Hadamard}{Paris}{Enseignant en mathématiques -- classes préparatoires scientifiques}
\begin{itemize}
  \item Préparer aux concours : cours, travaux dirigés et accompagnement individualisé en
  filières scientifiques.
\end{itemize}

\newpage

% ==================== PAGE 2 ====================
\begin{center}
  {\small\color{softgray}Louis Hauseux -- Percolia -- CV détaillé \textbullet{} 2/2}
\end{center}
\vspace{-0.75em}

\section*{Expériences professionnelles -- suite}
\cventry{2020--2022}{Activité concomitante}{Cours particuliers de mathématiques}{}{Professeur particulier}

\cventry{2020--2021}{Année scolaire}{Lycée Saint-Joseph}{Marvejols}{Professeur de mathématiques -- filières scientifiques}

\cventry{2013--2017}{4 ans}{École normale supérieure de Cachan}{Cachan / Paris}{Normalien élève -- fonctionnaire stagiaire}
\begin{itemize}
  \item Scolarité rémunérée après admission au concours d'informatique ; formation en
  mathématiques et informatique jusqu'à deux diplômes de Master 2.
\end{itemize}

% ==================== COMPÉTENCES ====================
\section*{Compétences}
\compactentry{Algorithmique / IA}{Clustering hiérarchique et non supervisé ; géométrie
algorithmique ; graphes, hypergraphes et complexes simpliciaux ; méthodes interprétables sans
apprentissage.}
\compactentry{3D / images}{Nuages de points LiDAR 3D/4D ; a priori géométriques ; traitement
d'images multimodales ; transport optimal et suivi temporel d'instances.}
\compactentry{Logiciel / calcul}{\textbf{Python, PyTorch, R, Git} ; calcul scientifique et GPU,
prototypage et optimisation, notebooks, traitement de grands jeux de données, \LaTeX.}
\compactentry{Transfert}{Cadrage de cas d'usage, benchmark, démonstrateur, communication
scientifique et collaboration avec chercheurs, ingénieurs et partenaires industriels.}

% ==================== FORMATION ====================
\section*{Formation}
\compactentry{2023--2026}{\textbf{Doctorat en informatique}, Université Côte d'Azur, École
doctorale STIC ; accueil Inria, financement DS4H/3IA ; direction K. Avrachenkov,
co-encadrement J. Zerubia. Soutenance prévue le 8 septembre 2026.}
\compactentry{2016--2017}{\textbf{Master 2 Probabilités et modèles aléatoires}, ENS Ulm /
École polytechnique / UPMC ; travail de recherche à Inria Paris sur le clustering par
complexes simpliciaux.}
\compactentry{2015--2016}{\textbf{Master 2 Mathématiques fondamentales}, ENS Ulm / École
polytechnique / UPMC ; premier diplôme de niveau master obtenu en 2016.}
\compactentry{2013--2015}{Licence de mathématiques et d'informatique, puis Master 1 de
mathématiques, ENS Cachan / ENS Ulm / UPMC / Université Paris-Diderot.}
\compactentry{2011--2013}{Classes préparatoires scientifiques MPSI--MP*, lycée Henri-IV,
Paris ; admission à l'ENS Cachan par le concours d'informatique.}

% ==================== PUBLICATIONS ====================
\section*{Publications et distinctions sélectionnées}
\compactentry{2026}{\textit{Generalization of Single-Linkage with Higher-Order Interactions},
avec K. Avrachenkov et J. Zerubia, \textit{Applied Network Science}.
\href{https://doi.org/10.1007/s41109-025-00756-1}{DOI}.}
\compactentry{2025}{\textit{How Can We Theoretically Measure the Performance of
Density-Based Clustering Algorithms?}, \textit{ACM SIGMETRICS Performance Evaluation
Review}. \href{https://doi.org/10.1145/3725536.3725546}{DOI}.}
\compactentry{2024--2025}{Publications sur les hypergraphes et le clustering (EUSIPCO,
Complex Networks) et sur l'extraction de fissures (GRETSI) ; liste complète sur
\href{https://inria.hal.science/search/index/?q=authIdHal_s:louis-hauseux}{HAL}.}
\compactentry{Distinctions}{Finaliste du prix Pierre Laffitte 2025 ; médaille d'excellence
de l'Université Côte d'Azur 2024 ; 2\textsuperscript{e} prix, catégorie \textit{Graduate},
ACM Student Research Competition SIGMETRICS 2024.}

\vfill
\begin{center}
  {\scriptsize\color{softgray}CV actualisé le 27 août 2026.}
\end{center}

\end{document}
